\chapter{‌متن برنامه‌ها}
\thispagestyle{empty}
در این پیوست، بخش‌های کلیدی از پیاده‌سازی معماری \lr{x2-DiT} ارائه شده است. تمرکز اصلی بر روی متد \lr{forward} است که در آن جریان ورودی با وضوح بالا (\lr{x2}) پردازش شده و توکن‌های کلاس محلی برای تجمیع ویژگی‌ها در آن درج می‌شوند. این کد نشان می‌دهد که چگونه دو جریان داده به‌صورت همزمان مدیریت می‌شوند.

\begin{latin}
\begin{python}
def forward(self, x, t, y):
    """
    Forward pass of x2-DiT.
    x: (N, C, H, W) tensor of spatial inputs
    t: (N,) tensor of diffusion timesteps
    y: (N,) tensor of class labels
    """

    # 1. Process high-resolution stream (x2)
    x2 = self.x2_embedder(x)  # (N, 4T, D)

    # Insert class tokens after every 4 patches to create local groups
    N, seq_len, D = x2.shape
    T = seq_len // 4

    # Reshape to group patches: (N, 4T, D) -> (N, T, 4, D)
    x2 = x2.reshape(N, T, 4, D)

    # Expand learnable class tokens: (1, T, D) -> (N, T, 1, D)
    x2_cls_tokens_expanded = self.x2_cls_tokens.repeat(N, 1, 1).unsqueeze(2)

    # Concatenate class tokens: (N, T, 4, D) + (N, T, 1, D) -> (N, T, 5, D)
    x2 = torch.cat([x2, x2_cls_tokens_expanded], dim=2)

    # Reshape back to sequence: (N, 5T, D)
    x2 = x2.reshape(N, 5 * T, D)

    # 2. Process main stream (standard DiT)
    x = self.x_embedder(x) + self.pos_embed  # (N, T, D)
    t = self.t_embedder(t)                   # (N, D)
    y = self.y_embedder(y, self.training)    # (N, D)
    c = t + y                                # (N, D)

    # The rest of the forward pass continues with block processing...
    # ...
\end{python}
\end{latin}

در کد فوق، مشاهده می‌شود که چگونه یک توکن کلاس یادگیرنده به ازای هر گروه ۴ تایی از پچ‌های \lr{x2} اضافه می‌شود. این ساختار به مدل اجازه می‌دهد تا اطلاعات محلی را به‌طور موثرتری نسبت به روش‌های سنتی مانند عملیات تجمیع و کاهش طول توالی ویژگی‌ها (\lr{Pooling}) جذب و نمایندگی کند.
